% Source: https://artofproblemsolving.com/community/c3h1788191_geometric_probability_problem
% Stable identity: aops-topic-1788191-q01
The probability is the area of the region of the rectangle between the circles $x^2+y^2=4$ and $x^2+y^2=15$, divided by the rectangle's area $15$.

\begin{asy}
usepackage("amsmath");
size(7.5cm);
defaultpen(fontsize(10pt));
import graph;
real r1=2, r2=sqrt(15);
pair O=(0,0), Q=(sqrt(6),3);
real th=degrees(atan(3/sqrt(6)));
path reg=arc(O,r2,0,th)--(0,3)--(0,r1)--arc(O,r1,90,0)--cycle;
fill(reg,lightgray);
draw((0,0)--(5,0)--(5,3)--(0,3)--cycle,linewidth(0.9));
draw(arc(O,r1,0,90),linewidth(0.8));
draw(arc(O,r2,0,th),linewidth(0.8));
draw(O--Q,dashed);
dot(Q);
label("$O$",O,SW);label("$(5,0)$",(5,0),S);label("$(5,3)$",(5,3),NE);label("$(0,3)$",(0,3),W);
label("$Q(\sqrt6,3)$",Q,N);
label("$2$",(2,0),S);label("$\sqrt{15}$",(sqrt(15),0),S);
draw(arc(O,0.75,0,th),linewidth(0.5));
label("$\theta$",1.05*dir(th/2));
label("\mbox{}",(min(currentpicture,user=true)+max(currentpicture,user=true))/2,black); // keeps the text after the figure black when LaTeX typesets the labels inline
\end{asy}

\textbf{Inside the larger circle.} Since $\sqrt{15}<5$, the circle of radius $\sqrt{15}$ does not reach the side $x=5$, but it crosses the top side $y=3$ at $Q=(\sqrt6,3)$, because $6+9=15$. The part of the rectangle inside this circle is the right triangle with vertices $O$, $(0,3)$, $Q$, of area $\dfrac{1}{2}\cdot3\cdot\sqrt6=\dfrac{3\sqrt6}{2}$, together with the circular sector between the $x$-axis and $OQ$, of angle $\theta=\arctan\dfrac{3}{\sqrt6}=\arctan\dfrac{\sqrt6}{2}$ and area $\dfrac{1}{2}\cdot15\,\theta$.

\textbf{Inside the smaller circle.} The quarter disc of radius $2$ lies inside the rectangle (as $2<3$) and has area $\dfrac{1}{4}\cdot4\pi=\pi$.

So the shaded region has area $\dfrac{3\sqrt6}{2}+\dfrac{15}{2}\arctan\dfrac{\sqrt6}{2}-\pi$, and the probability is
\[
\boxed{\dfrac{\sqrt6}{10}+\dfrac{1}{2}\arctan\dfrac{\sqrt6}{2}-\dfrac{\pi}{15}}\approx0.479 .
\]
